\chapter{Retournement de court terme}\label{ch:s1:short-term-reversal}

Acheter les perdants de la veille et vendre les gagnants de la veille. Sur le marché synthétique, la version résiduelle de cette règle présente un coefficient d'information de rang (rank IC) de 0.041
face au rendement du lendemain, et le portefeuille naïf construit à partir de celui-ci (exposition brute égale à un, reconstruit à chaque clôture, sur les 500 valeurs les plus liquides) affiche un ratio de Sharpe de
13.0 avant coûts. À dix points de base par unité négociée, il tombe à $-3.4$. La stratégie est réelle (c'est l'une des plus anciennes anomalies de la littérature, et le
marché synthétique l'introduit) et quasiment inexploitable telle quelle, car elle fait tourner les trois quarts de son portefeuille chaque jour. Ce chapitre porte sur le retournement à un jour
et de ses variantes, et surtout de la manière dont on la négocie~: en intégrant les coûts dans l'optimiseur, elle rapporte un ratio de Sharpe de 1.7 après coûts pour un encours de 100~millions de dollars
et à peu près rien à 1~milliard de dollars. La construction est \texttt{firm.reversal}.

\section{Pourquoi les prix se retournent~: payer pour l'immédiateté}

\begin{definition}[Retournement de court terme]\label{def:s1:short-term-reversal:reversal}
Le \emph{retournement de court terme}\index{retournement de court terme} est la tendance des actions ayant de faibles rendements sur une courte fenêtre passée (un jour, une semaine, un mois) à
surperformer, sur la fenêtre suivante de même durée, les actions ayant de forts rendements~; une stratégie de retournement de court terme achète les perdants récents et vend les gagnants récents.
\end{definition}

L'effet est présent dans la littérature depuis aussi longtemps qu'il existe des fichiers de rendements. Jegadeesh a mis en évidence une autocorrélation sérielle d'ordre~1 négative et significative dans
les rendements mensuels, les déciles extrêmes d'une prévision fondée sur elle différant de 2.49~\% par mois sur 1934--1987~; Lehmann constate que les gagnants et les perdants d'une semaine
s'inversent nettement la semaine suivante, avec des profits qui résistent à la correction des fourchettes bid--ask et à des coûts plausibles. La lecture de Lehmann est devenue la référence~:
le retournement est «~une inefficience du marché de la liquidité autour des fortes variations de prix~». Un trader qui doit vendre tout de suite fait baisser le cours, qui
remonte à mesure que d'autres absorbent le flux, et celui qui a acheté au vendeur pressé est rémunéré pour avoir fourni de l'immédiateté. Nagel a rendu cette interprétation explicite, en traitant
les rendements de retournement comme approximation des rendements tirés de la fourniture de liquidité, et il les trouve maximaux quand la liquidité est la plus rare~: les rendements espérés et les
ratios de Sharpe conditionnels bondissent en période de turbulences financières, de façon prévisible avec le VIX.

Deux conséquences structurent tout ce qui suit. D'abord, la part d'un mouvement de prix qui relève de l'information ne devrait pas s'inverser, et la part qui relève de la pression devrait s'inverser, si bien qu'un bon
signal de retournement cherche à isoler la pression. Ensuite, le trader en retournement est un teneur de marché à horizon quotidien~: l'edge est la rémunération de la négociation, et la
stratégie tient ou tombe selon le coût de cette négociation.

\section{Retournement quotidien~: brut, résiduel, ajusté du secteur}

\begin{definition}[Retournement résiduel, retournement ajusté du secteur]\label{def:s1:short-term-reversal:residual}
Le \emph{retournement résiduel}\index{retournement résiduel} est le \omterm{def:s1:short-term-reversal:reversal}{retournement de court terme} calculé sur les rendements résiduels, soit la part du rendement de chaque action que ses
expositions aux facteurs n'expliquent pas. Le \emph{retournement ajusté du secteur}\index{retournement ajusté du secteur} est le retournement calculé sur le rendement de chaque action en excès du
rendement moyen de son secteur.
\end{definition}

La journée d'une action est celle du marché, celle de son secteur et la sienne. Les composantes de marché et de secteur, partagées par des centaines d'actions, reflètent pour l'essentiel
l'information sur l'économie~; ce ne sont pas elles qu'un vendeur pressé déplace. Le marché synthétique (livre~7, chapitre~5) introduit un retournement de 5~\% du choc spécifique de chaque jour et
rien sur les composantes communes, de sorte que plus l'estimation du choc spécifique par le signal est nette, mieux il prévoit~:

\begin{center}
\small
\begin{tabular}{@{}lcccc@{}}
\toprule
années 3 à 10 & brut & secteur & résiduel & vérité\\
\midrule
IC de rang moyen, toutes les valeurs cotées & 0.037 & 0.038 & 0.041 & 0.044\\
IC de rang moyen, 500 valeurs les plus liquides & 0.037 & 0.039 & 0.041 & 0.044\\
\midrule
portefeuille naïf~: ratio de Sharpe avant coûts & 6.69 & 7.24 & 12.96 & \\
rendement annuel~; rotation quotidienne & 27.0~\%~; 75~\% & 28.2~\%~; 75~\% & 30.3~\%~; 76~\% & \\
ratio de Sharpe après 10 pb & $-2.62$ & $-2.46$ & $-3.43$ & \\
coût de point mort par unité négociée & 7.2 pb & 7.5 pb & 7.9 pb & \\
\bottomrule
\end{tabular}
\end{center}

(Les portefeuilles naïfs négocient les 500 valeurs les plus liquides~; «~secteur~» désigne le signal ajusté du secteur et «~vérité~» le
rendement espéré injecté.) Le coefficient d'information ne progresse que de 0.037 à 0.041, mais le ratio de Sharpe du portefeuille résiduel presque double, parce que le portefeuille brut porte
des expositions au secteur et au style qui ajoutent du risque sans rendement. Blitz, Huij, Lansdorp et Verbeek ont observé la même chose sur des données américaines~: les stratégies de retournement classiques ont des
expositions dynamiques aux facteurs de Fama--French, et la version résiduelle, sans elles, procure des rendements ajustés du risque deux fois plus élevés, significatifs nets des coûts de transaction
même parmi les grandes capitalisations après 1990. Da, Liu et Schaumburg sont allés plus loin dans la même direction, en retirant la part du rendement passé expliquée par l'information sur les flux de trésorerie
(mesurée au moyen des révisions des prévisions des analystes)~; le rendement ajusté du risque de leur stratégie améliorée valait quatre fois celui de la version standard.

Aucun de ces chiffres bruts ne survit aux coûts sous forme de portefeuille naïf. Chaque portefeuille du tableau fait tourner environ les trois quarts de son exposition brute chaque jour~; à dix points de base
par unité de poids négociée, il coûte 38~\% par an, contre un rendement brut de 27--30~\%. Le coût de point mort, égal au rendement brut divisé par deux fois la rotation (livre~7, chapitre~16),
se situe entre 7 et 8 points de base pour les trois, et dix points de base est déjà un coût faible pour un portefeuille quotidien (le livre~7, chapitre~16 faisait le même constat avec le signal brut).

\section{Retournement intrajournalier et retournement overnight}

\begin{definition}[Retournement overnight]\label{def:s1:short-term-reversal:overnight}
Le \emph{retournement overnight}\index{retournement overnight} est la tendance du rendement overnight d'une action, de la clôture précédente à l'ouverture, à être en partie inversé
pendant la séance suivante, et celle du rendement intrajournalier à être en partie inversé pendant la nuit~; une stratégie fondée sur elle négocie entre les deux moitiés de la
journée.
\end{definition}

Découper la journée en deux change le tableau. Lou, Polk et Skouras ont mis en évidence une forte continuation au sein de chaque partie (les rendements overnight prédisent les rendements overnight, et
les rendements intrajournaliers prédisent les rendements intrajournaliers pendant des années), ainsi qu'un retournement inter-périodes qui les compense, et dans quatorze stratégies, les profits ont été réalisés soit
entièrement overnight ou entièrement intrajournaliers~; pour le retournement, overnight. Akbas, Boehmer, Jiang et Koch ont compté, chaque mois, la fréquence à laquelle un rendement overnight positif
était suivi d'un retournement négatif en séance, et ont constaté que les actions présentant davantage de tels jours (un «~bras de fer~» plus intense entre les clientèles
overnight et de jour) rapportaient ensuite des rendements plus élevés. À l'intérieur de la journée, Heston, Korajczyk et Sadka ont relié le \omterm{def:s1:short-term-reversal:reversal}{retournement de court terme} à des déséquilibres temporaires de liquidité d'une durée
inférieure à une heure et au rebond bid--ask, à côté d'une remarquable continuation à des intervalles qui sont des multiples exacts d'un jour.

\texttt{firm.reversal} calcule les signaux overnight et intrajournalier à partir des cours d'ouverture et de clôture (fonctions \texttt{overnight} et \texttt{intraday})~;
le marché synthétique n'a pas d'ouverture, si bien que le chapitre ne mesure que les variantes quotidiennes, et les deux fiches de stratégie qui s'y rapportent reposent sur les résultats publiés.

\section{Coûts~: là où la stratégie tient ou tombe}

Un portefeuille de retournement quotidien a besoin de trois choses que le portefeuille naïf n'a pas. Une prévision exprimée en unités de rendement, pour pouvoir la comparer aux coûts~: la règle de Grinold
(livre~7, chapitre~15) donne $\alpha_i = \mathrm{IC}\,\sigma_i\,z_i$, avec $z_i$ le résiduel rapporté à la volatilité spécifique propre de l'action et standardisé entre
les valeurs, et l'IC de 0.041 mesuré sur la seule année~2. Un modèle de coûts~: une demi-fourchette de 2 points de base plus l'impact en racine carrée
$0.7\,\sigma_i\sqrt{Q_i/V_i}$ (livre~7, chapitre~27). Enfin, un optimiseur qui compare les deux pour chaque valeur et chaque jour. Avec un modèle de risque diagonal, le problème se découpe par valeur~:
le trade de chaque valeur est un pas en forme fermée qui réduit le trade souhaité du coût linéaire et l'amortit par l'impact, et la neutralité aux
quinze facteurs du modèle de risque ajoute un petit système pour les multiplicateurs (\cref{lst:s1:short-term-reversal:book}).

\omterm{def:s1:short-term-reversal:residual}{Retournement résiduel} sur toutes les valeurs, jours de publication des résultats exclus~:

\begin{center}
\small
\begin{tabular}{@{}p{6.6cm}ccc@{}}
\toprule
taille du fonds & 100~M\$ & 1~Md\$ & 10~Md\$\\
\midrule
coûts ignorés~: ratio de Sharpe avant, après & 16.92, $-27.8$ & 16.92, $-31.7$ & 16.92, $-32.2$\\
coûts ignorés~: coûts de transaction par an & 246~\% & 738~\% & 2\,293~\%\\
coûts intégrés~: ratio de Sharpe avant coûts & 8.42 & 3.97 & 1.72\\
coûts intégrés~: ratio de Sharpe après coûts & 1.72 & 0.09 & 0.19\\
coûts intégrés~: brut, coûts, net (en~\% par an) & 10.23, 8.18, 2.05 & 2.44, 2.38, 0.05 & 0.45, 0.40, 0.05\\
coûts intégrés~: rotation quotidienne~; brut & 18.6~\%~; 0.91 & 3.9~\%~; 0.48 & 0.5~\%~; 0.18\\
profit net annuel & 2.0~M\$ & 0.5~M\$ & 4.8~M\$\\
\bottomrule
\end{tabular}
\end{center}

Le livre sans coûts traite 186~\% de son capital par jour, pour une exposition brute de 2.48~: son ratio de Sharpe brut de 16.9 est le meilleur de tous les livres du chapitre, et
à 100 millions de dollars il paie 246~\% de son capital par an en coûts. Coûts intégrés, le livre à 100 millions de dollars traite dix fois moins et conserve la moitié de son
ratio de Sharpe brut, et dégage un ratio net de 1.72. À 1 milliard de dollars, il traite 3.9~\% par jour et ne fait que couvrir ses coûts~; à 10 milliards de dollars, il ne traite presque plus
(\cref{fig:s1:short-term-reversal:size}). L'edge par transaction de la stratégie est son coefficient d'information multiplié par la volatilité d'une journée, soit environ seize points de base
pour un titre situé à deux écarts-types, et un coût en racine carrée atteint ce niveau sur une transaction représentant 1~\% du volume quotidien du titre. La capacité d'un retournement à un jour
est donc faible~: pour ce signal, elle se situe entre 100 millions et 1 milliard de dollars.

\begin{omfigure}
\begin{tikzpicture}
\begin{axis}[width=8.4cm,height=5.4cm,xlabel={\small taille du fonds},ylabel={\small ratio de Sharpe},tick label style={font=\footnotesize},
  xtick={8,8.477,9,9.477,10},xticklabels={100~M\$,300~M\$,1~Md\$,3~Md\$,10~Md\$},ymin=-1,ymax=9,grid=major,grid style={gray!25},
  legend pos=north east,legend style={font=\scriptsize},table/col sep=comma]
\addplot[omDef,very thick,mark=*] table[x=log10_aum,y=sr_gross]{figdata/strategies-1/02-short-term-reversal/size.csv};
\addplot[omThm,very thick,mark=square*] table[x=log10_aum,y=sr_net]{figdata/strategies-1/02-short-term-reversal/size.csv};
\addplot[black!50,dashed] coordinates {(8,0) (10,0)};
\legend{avant coûts,après coûts}
\end{axis}
\end{tikzpicture}

\omcaption{\omterm{def:s1:short-term-reversal:residual}{Retournement résiduel} avec coûts intégrés dans l'optimiseur, sur le marché synthétique~: ratio de Sharpe avant et après coûts en fonction de la taille du fonds. Données~:
\texttt{s1\_reversal.traded}.}\label{fig:s1:short-term-reversal:size}
\end{omfigure}

Un filtre supplémentaire a fait la différence entre une petite perte et un petit profit. Le marché synthétique intègre une dérive post-annonce (livre~7, chapitre~11), et un
titre qui chute à la suite de mauvais résultats continue de dériver à la baisse~: un livre de retournement qui l'achète parie contre l'information. Ignorer les titres le jour de leur annonce, qui
est programmée et publique, coûte presque rien en coefficient d'information (0.040 contre 0.041) et fait passer le livre à 1 milliard de dollars de $-0.70~\%$ par an
(ratio de Sharpe $-1.04$) à $+0.05\%$ (\cref{fig:s1:short-term-reversal:earnings}). C'est le propos de Da, Liu et Schaumburg, à petite échelle~: les mouvements expliqués par
les fondamentaux ne se retournent pas.

\begin{omfigure}
\begin{tikzpicture}
\begin{axis}[width=9cm,height=5.4cm,xlabel={\small année},ylabel={\small P\&L net cumulé (\% du capital)},tick label style={font=\footnotesize},
  xmin=2,xmax=10,ymin=-7,ymax=4,grid=major,grid style={gray!25},legend pos=north east,legend style={font=\scriptsize},table/col sep=comma]
\addplot[omDef,very thick] table[x=year,y=filtered_pct]{figdata/strategies-1/02-short-term-reversal/earnings.csv};
\addplot[omThm,thick] table[x=year,y=unfiltered_pct]{figdata/strategies-1/02-short-term-reversal/earnings.csv};
\legend{jours d'annonce écartés,chaque jour}
\end{axis}
\end{tikzpicture}

\omcaption{Le livre de retournement sensible aux coûts à 1 milliard de dollars, net de coûts, avec et sans le filtre sur les résultats. Données~: \texttt{s1\_reversal.traded}.}\label{fig:s1:short-term-reversal:earnings}
\end{omfigure}

\section{Le déclin des profits du retournement}

Les données publiques confirment les données synthétiques sur les coûts et y ajoutent une tendance. Le facteur de \omterm{def:s1:short-term-reversal:reversal}{retournement de court terme} de Kenneth French (pondéré par la valeur, écart entre les rendements du mois précédent les plus
faibles et les plus élevés, grandes et petites capitalisations moyennées, rééquilibré mensuellement) a rapporté 10.9~\% par an de 1926 à 1989, avec un ratio de Sharpe de 0.93 ($t = 7.42$), et
1.6~\% par an de 1990 à juillet 2026, avec un ratio de Sharpe de 0.13 ($t = 0.79$), avant tout coût de transaction~; depuis le début des années 2020, il a perdu 5.2~\% par an
(\cref{fig:s1:short-term-reversal:decline}). Il s'agit d'un retournement mensuel sur de grands portefeuilles pondérés par la valeur, non de la version résiduelle quotidienne, et il n'y a pas là
de preuve que l'effet quotidien a disparu~; les résultats de Nagel indiquent que ses rendements se concentrent en période de crise. Mais le long déclin du facteur classique,
quelles qu'en soient les causes (la publication, la baisse des coûts de négociation et l'afflux de capitaux en concurrence pour fournir la liquidité sont les candidats habituels, et ces données ne permettent pas
de les départager), constitue le contexte de tout livre de retournement. L'edge qui subsiste revient à
celui qui le négocie au moindre coût.

\begin{omfigure}
\begin{tikzpicture}
\begin{axis}[width=10cm,height=5.0cm,xlabel={\small année},ylabel={\small moyenne sur 10 ans (\% par an)},tick label style={font=\footnotesize},
  xmin=1935,xmax=2026,ymin=-5,ymax=30,grid=major,grid style={gray!25},table/col sep=comma,
  x tick label style={/pgf/number format/1000 sep={}}]
\addplot[omDef,very thick] table[x=year,y=mean_pct]{figdata/strategies-1/02-short-term-reversal/rolling.csv};
\addplot[black!50,dashed] coordinates {(1935,0) (2026,0)};
\addplot[omThm,thick,dashed] coordinates {(1990,-5) (1990,30)};
\node[font=\scriptsize,text=omThm,anchor=west] at (axis cs:1991,24) {publié en 1990};
\end{axis}
\end{tikzpicture}

\omcaption{Facteur de \omterm{def:s1:short-term-reversal:reversal}{retournement de court terme} de Kenneth French~: moyenne glissante sur 120 mois des rendements, annualisée, à chaque mois de décembre~; la ligne en pointillés marque
les articles de Jegadeesh et Lehmann. Dérivé de la Kenneth R. French
Data Library (ST\_Rev mensuel, fichier CRSP 202607)~; la série brute n'est pas redistribuée.}\label{fig:s1:short-term-reversal:decline}
\end{omfigure}

\section{Fiches de stratégie}

\begin{strategyfile}{Retournement résiduel quotidien}\label{strat:s1:short-term-reversal:daily}
\sfield{Qui paie, et pourquoi} Les traders qui doivent vendre ou acheter tout de suite, et qui font bouger les cours pour y parvenir~; le livre est rémunéré pour absorber leurs flux et attendre un jour.
\sfield{Instruments et plateformes} Actions individuelles liquides, en position longue et courte, sur leurs bourses principales~; les transactions se croisent à la clôture ou pendant la séance suivante.
\sfield{Signal} Opposé du rendement résiduel du jour sur un modèle factoriel de risque, normalisé par la volatilité spécifique de l'action~; on écarte les titres ayant des résultats ou d'autres annonces
programmées ce jour-là.
\sfield{Dimensionnement et exécution} Une prévision selon la règle de Grinold, un optimiseur sensible aux coûts et neutre vis-à-vis des facteurs du modèle, des transactions exécutées jusqu'à la clôture.
\sfield{Coûts} L'ensemble de la stratégie~: un taux de rotation de plusieurs dizaines de pourcents par jour~; une demi-fourchette et un impact sur chaque transaction~; l'emprunt sur les positions courtes.
\sfield{Comment elle meurt} Les coûts augmentent ou le livre dépasse sa capacité~; les jours d'annonce contaminent le signal~; d'autres fournisseurs de liquidité font disparaître l'edge.
\sfield{Horizon, capacité, infrastructure} D'un à quelques jours~; capacité faible (pour le signal du chapitre, en dessous de 1 milliard de dollars)~; un modèle de risque quotidien, un modèle de coûts
calibré sur les exécutions de la firme, un système d'exécution aux enchères de clôture.
\sfield{Backtest honnête} Résiduels et univers tels qu'observés à chaque date~; un modèle de coûts calibré sur des exécutions réelles, appliqué à chaque transaction~; le filtre sur les résultats, connu à l'avance~;
résultats par taille.
\sfield{Sources} Blitz, Huij, Lansdorp et Verbeek (2013)~; Da, Liu et Schaumburg (2014)~; la simulation de ce chapitre.
\end{strategyfile}

\begin{strategyfile}{Retournement intrajournalier à la clôture}\label{strat:s1:short-term-reversal:intraday}
\sfield{Qui paie, et pourquoi} Les traders qui poussent les cours pendant la séance avec des déséquilibres temporaires~; le livre fournit la contrepartie et se déboucle à mesure que le
déséquilibre se résorbe.
\sfield{Instruments et plateformes} Actions liquides et leurs plateformes lit et dark pendant la séance.
\sfield{Signal} Opposé du rendement de l'action sur les dernières minutes, jusqu'à une heure, relatif à son secteur ou au marché, avec des conditions de volume et de déséquilibre.
\sfield{Dimensionnement et exécution} Ordres passifs lorsque c'est possible~; positions clôturées avant la fin de la journée.
\sfield{Coûts} La fourchette et la sélection adverse dominent~; l'edge est de l'ordre d'une fourchette.
\sfield{Comment elle meurt} Des fournisseurs de liquidité plus rapides prennent d'abord le déséquilibre~; le mouvement s'avère être une information~; le rebond bid-ask pris pour un retournement.
\sfield{Horizon, capacité, infrastructure} De quelques minutes à quelques heures~; capacité faible~; données intrajournalières, un simulateur de carnet d'ordres (livre~7, chapitre~18) et une exécution à faible
latence.
\sfield{Backtest honnête} Prix de transaction et non prix milieux dans le signal, et prix milieux dans l'évaluation~; un modèle d'exécution tenant compte de la file d'attente.
\sfield{Sources} Heston, Korajczyk et Sadka (2010).
\end{strategyfile}

\begin{strategyfile}{Retournement overnight}\label{strat:s1:short-term-reversal:overnight}
\sfield{Qui paie, et pourquoi} Les deux clientèles du bras de fer~: la demande de l'ouverture et de la nuit, qui pousse les cours, et les traders de la journée, qui les repoussent.
\sfield{Instruments et plateformes} Actions liquides~; les enchères d'ouverture et de clôture.
\sfield{Signal} Opposé du rendement overnight, négocié pendant la séance~; ou opposé du rendement de la séance, négocié la nuit.
\sfield{Dimensionnement et exécution} Entrer à l'enchère d'ouverture et sortir à la clôture, ou l'inverse~; neutre vis-à-vis du marché et des secteurs.
\sfield{Coûts} Deux croisements aux enchères par jour~; les prix d'enchère peuvent évoluer contre les gros ordres.
\sfield{Comment elle meurt} La continuation à l'intérieur de chaque partie de la journée joue contre lui~; la configuration change selon les acteurs qui négocient à l'ouverture.
\sfield{Horizon, capacité, infrastructure} Une demi-journée~; capacité limitée par le volume des enchères~; prix d'ouverture et de clôture, données de déséquilibre d'enchère.
\sfield{Backtest honnête} Prix officiels d'ouverture et de clôture, et non les premières et dernières transactions~; coûts d'enchère~; aucun signal issu d'une ouverture à laquelle on ne pourrait pas trader.
\sfield{Sources} Lou, Polk et Skouras (2019)~; Akbas, Boehmer, Jiang et Koch (2022).
\end{strategyfile}

\begin{strategyfile}{Retournement en surcouche de fourniture de liquidité}\label{strat:s1:short-term-reversal:overlay}
\sfield{Qui paie, et pourquoi} Les mêmes traders pressés, mais le retournement n'est pas négocié seul~: il infléchit les transactions qu'un portefeuille plus lent réalise de toute façon.
\sfield{Instruments et plateformes de négociation} Les portefeuilles d'actions existants de la maison.
\sfield{Signal} La prévision de retournement ajoutée à la prévision du portefeuille plus lent dans l'optimiseur, de sorte que le portefeuille achète ce qu'il avait prévu d'acheter les jours de baisse et vende les
jours de hausse.
\sfield{Dimensionnement et exécution} Par l'optimiseur du portefeuille plus lent~; les transactions sont celles qu'il réaliserait de toute façon, avec un calendrier décalé.
\sfield{Coûts} Presque aucun coût propre~: la surcouche modifie le moment où les transactions sont réalisées, et non leur nombre.
\sfield{Comment elle meurt} Elle ne meurt pas seule~; elle s'affaiblit lorsque le portefeuille plus lent négocie moins.
\sfield{Horizon, capacité, infrastructure} Quelques jours~; capacité fixée par la rotation du portefeuille hôte~; l'optimiseur du portefeuille hôte.
\sfield{Backtest honnête} Comparer le portefeuille hôte avec et sans la surcouche, à même rotation et mêmes coûts.
\sfield{Sources} Nagel (2012), sur les rendements de retournement comme rémunération de la fourniture de liquidité~; livre~7, chapitre~27.
\end{strategyfile}

\section{Tutoriel~: sept avant les coûts}\label{tut:s1:short-term-reversal:tut}

\begin{tutorial}
\textbf{Objectif.} Mesurer les trois retournements quotidiens, puis négocier le \omterm{def:s1:short-term-reversal:residual}{retournement résiduel} en intégrant les coûts dans l'optimiseur, pour trois tailles de fonds. \textbf{État final~:} les deux
tableaux et la \cref{fig:s1:short-term-reversal:size}.
\begin{tutsteps}
\item \textbf{Signaux et IC}~: \texttt{s1\_reversal.signals}, \texttt{ic} pour les signaux bruts, ajustés du secteur et résiduels, et pour les rendements réalisés, avec et sans jours d'annonces de résultats.
\item \textbf{Les portefeuilles naïfs}~: \texttt{naive(kind)} et \texttt{naive(kind, 0.001)}, et \texttt{breakeven}.
\item \textbf{Le portefeuille intégrant les coûts.} Un pas en forme close par action et la méthode de Newton pour les multiplicateurs de neutralité aux facteurs.
\omcode{code/firm/reversal/firm_reversal.py}{69}{97}{Le portefeuille du jour, intégrant les coûts et neutre aux facteurs.}\label{lst:s1:short-term-reversal:book}
\item \textbf{Chaque jour}~: prévoir, optimiser, imputer le coût, laisser dériver le portefeuille.
\omcode{code/strategies-1/02-short-term-reversal/python/s1_reversal.py}{137}{151}{La boucle quotidienne du portefeuille négocié.}\label{lst:s1:short-term-reversal:loop}
\item \textbf{Lancer} \texttt{traded(aum, costs, 1.0, skip\_earnings)} pour 100~millions de dollars, 1~milliard de dollars et 10~milliards de dollars, et \texttt{fig\_reversal.py}.
\end{tutsteps}
\textbf{Pour aller plus loin.} Lisser la prévision (livre~7, chapitre~27) et voir si un portefeuille plus lent bat un portefeuille plus rapide~; restreindre l'univers aux 500 plus
liquides~; multiplier par deux le coût dans l'optimiseur, comme si chaque transaction devait être débouclée, et expliquer pourquoi cela n'aide pas ici.
\end{tutorial}

\section{Construction~: la famille des retournements}\label{bld:s1:short-term-reversal:reversal}

\begin{build}
\bfield{Rôle} Les signaux de la famille des retournements et un portefeuille quotidien qui les négocie seulement là où ils battent leurs coûts.
\bfield{Interface} \texttt{raw(r)}, \texttt{industry\_adjusted(r, industry, listed)}, \texttt{residual(r, X, w, listed)},
\texttt{overnight(open\_, prev\_close)}, \texttt{intraday(close, open\_)}, \texttt{forecast(z, ic, sigma)},
\texttt{book(alpha, var, w0, gamma, aum, sigma, adv, half\_spread, eta, costs, round\_trip, X)}.
\bfield{Règles} Signaux issus des seules données du jour~; prévisions exprimées en unités de rendement~; le coût dans l'optimiseur est le coût imputé~; neutralité exacte aux
expositions détenues.
\bfield{Tests d'acceptation} \texttt{code/firm/reversal/tests/}~: les signaux vérifiés à la main~; le résidu orthogonal aux expositions~; sans coûts, le portefeuille est égal à la
solution moyenne--variance en forme close~; les coûts réduisent les transactions~; la neutralité aux facteurs est respectée à $10^{-10}$ près avec les coûts.
\bfield{Pour aller plus loin} Une matrice de covariance factorielle complète dans l'optimiseur (\texttt{cost\_aware} du livre~7, chapitre~27)~; négociation multipériode vers un portefeuille cible (livre~7,
chapitre~26)~; les signaux overnight et intrajournaliers sur de vrais cours d'ouverture et de clôture.
\end{build}

\begin{omsources}
\item N.~Jegadeesh, ``Evidence of predictable behavior of security returns'', \emph{Journal of Finance} 45(3), 1990.
\item B.~N.~Lehmann, ``Fads, martingales, and market efficiency'', \emph{Quarterly Journal of Economics} 105(1), 1990.
\item S.~Nagel, ``Evaporating liquidity'', \emph{Review of Financial Studies} 25(7), 2012.
\item D.~Blitz, J.~Huij, S.~Lansdorp and M.~Verbeek, ``Short-term residual reversal'', \emph{Journal of Financial Markets} 16(3), 2013.
\item Z.~Da, Q.~Liu and E.~Schaumburg, ``A closer look at the short-term return reversal'', \emph{Management Science} 60(3), 2014.
\item S.~L.~Heston, R.~A.~Korajczyk and R.~Sadka, ``Intraday patterns in the cross-section of stock returns'', \emph{Journal of Finance} 65(4), 2010.
\item D.~Lou, C.~Polk and S.~Skouras, ``A tug of war: overnight versus intraday expected returns'', \emph{Journal of Financial Economics} 134(1), 2019.
\item F.~Akbas, E.~Boehmer, C.~Jiang and P.~D.~Koch, ``Overnight returns, daytime reversals, and future stock returns'', \emph{Journal of Financial Economics}
145(3), 2022.
\item Kenneth R. French Data Library, Short-Term Reversal Factor (ST Rev), mensuel.
\end{omsources}

\section{Exercices}

\begin{exercise}[$\star$]\label{exo:s1:short-term-reversal:1}
Le portefeuille naïf résiduel rapporte 30.3~\% par an avant coûts et fait tourner 76~\% de son exposition brute chaque jour. Calculer le coût de point mort par unité négociée.
\end{exercise}

\begin{exercise}[$\star$]\label{exo:s1:short-term-reversal:2}
Une action dont le prix milieu ne bouge pas pendant deux jours~; la fourchette vaut 10 points de base. La dernière transaction du premier jour a lieu au cours acheteur, et celle du second jour au cours vendeur. Quel rendement
un fichier close-to-close enregistre-t-il pour le second jour, et que conclut un signal de retournement fondé sur les prix de transaction~?
\end{exercise}

\begin{exercise}[$\star$]\label{exo:s1:short-term-reversal:3}
Avec un impact en racine carrée $0.7\,\sigma\sqrt{Q/V}$, quel est l'impact d'une transaction de 10~millions de dollars dans une action de volatilité quotidienne de 2~\% et de 25~millions de dollars de volume quotidien~?
\end{exercise}

\begin{exercise}[$\star\star$]\label{exo:s1:short-term-reversal:4}
L'IC vaut 0.04 et la volatilité quotidienne d'une action est de 2~\%. Quelle est la prévision de Grinold pour un titre situé à deux écarts-types, et quelle taille de transaction sur l'action
de l'exercice~3 peut-elle couvrir, avec une demi-fourchette de 2 points de base~?
\end{exercise}

\begin{exercise}[$\star\star$]\label{exo:s1:short-term-reversal:5}
Pourquoi le portefeuille de \omterm{def:s1:short-term-reversal:residual}{retournement résiduel} a-t-il presque deux fois le ratio de Sharpe du portefeuille brut alors que son coefficient d'information n'est supérieur que de 10~\%~?
\end{exercise}

\begin{exercise}[$\star\star$]\label{exo:s1:short-term-reversal:6}
Expliquer pourquoi un portefeuille de retournement qui achète des actions le jour d'une mauvaise annonce de résultats perd de l'argent sur un marché avec dérive post-annonce, et pourquoi ignorer
les jours d'annonce n'introduit pas de biais de look-ahead.
\end{exercise}

\begin{exercise}[$\star\star\star$]\label{exo:s1:short-term-reversal:7}
\emph{Programmation.} Lancer \texttt{traded(3e8, True, 1.0, True)}. Indiquer le ratio de Sharpe net, et expliquer pourquoi le ratio de Sharpe net n'est pas une fonction monotone de la taille du
fonds dans la \cref{fig:s1:short-term-reversal:size}.
\end{exercise}

\begin{exercise}[$\star\star\star$]\label{exo:s1:short-term-reversal:8}
\emph{Trouver la faille.} «~Notre backtest de \omterm{def:s1:short-term-reversal:residual}{retournement résiduel} affiche un ratio de Sharpe de 13. Même si la négociation réelle le divise par deux, une allocation de 1~milliard de dollars est sûre.~»
\end{exercise}

\section{Problème~: sept avant les coûts}

\begin{problem}[{Problème du week-end --- une stratégie qui vit de ses coûts}]\label{pb:s1:short-term-reversal:1}
Les signaux de retournement du chapitre sur \texttt{firm.synthmkt}, et le facteur de retournement sur des données réelles.

\medskip\noindent\textbf{Partie~I --- L'effet.}
\begin{enumerate}
\item Définir le \omterm{def:s1:short-term-reversal:reversal}{retournement de court terme} et donner l'explication par la fourniture de liquidité.
\item Que mettent en évidence Jegadeesh et Lehmann~?
\item Que montre Nagel sur les périodes où le retournement rapporte~?
\item Pourquoi les nouvelles ne devraient-elles pas se retourner~?
\end{enumerate}

\medskip\noindent\textbf{Partie~II --- Les signaux.}
\begin{enumerate}[resume]
\item Donner les IC du retournement brut, ajusté du secteur et résiduel, ainsi que celui de la vérité.
\item Donner les ratios de Sharpe des portefeuilles naïfs avant et après dix points de base.
\item Pourquoi le portefeuille résiduel fait-il nettement mieux avant coûts~?
\item Que trouvent Blitz et ses coauteurs, ainsi que Da, Liu et Schaumburg~?
\end{enumerate}

\medskip\noindent\textbf{Partie~III --- Mise en œuvre.}
\begin{enumerate}[resume]
\item Comment la prévision est-elle exprimée en unités de rendement~?
\item Décrire l'optimiseur et expliquer pourquoi il traite séparément chaque titre.
\item Que négocie le portefeuille aveugle aux coûts, et combien cela lui coûte-t-il~?
\item Que fait le filtre sur les résultats, et que vaut-il à 1~milliard de dollars~?
\end{enumerate}

\medskip\noindent\textbf{Partie~IV --- Le verdict.}
\begin{enumerate}[resume]
\item Énoncer le \emph{résultat nommé}~: le ratio de Sharpe net du \omterm{def:s1:short-term-reversal:residual}{retournement résiduel} à 100~millions de dollars et à 1~milliard de dollars, coûts inclus dans l'optimiseur, face
à celui du portefeuille naïf.
\item Où se situe la capacité de la stratégie, et pourquoi est-elle si faible~?
\item Que s'est-il passé pour le facteur de retournement de French après 1989~?
\item Pourquoi cela ne prouve-t-il pas que l'effet quotidien a disparu~?
\item Qu'apportent les décompositions en rendements overnight et en rendements intrajournaliers~?
\item Quelle fiche de stratégie lancer avec 50~millions de dollars, et laquelle avec 5~milliards de dollars~?
\item Quelle est la donnée d'entrée la plus importante d'un backtest de retournement~?
\item En une phrase~: qu'est-ce que le \omterm{def:s1:short-term-reversal:reversal}{retournement de court terme}, d'un point de vue économique~?
\end{enumerate}
\end{problem}

\section{Questions d'entretien}

\begin{interviewq}[$\star$ \iqroles{chercheur, trader}]\label{iq:s1:short-term-reversal:1}
Pourquoi les rendements de court terme se retournent-ils, et qui paie le trader de retournement~?
\end{interviewq}

\begin{interviewq}[$\star\star$ \iqroles{chercheur}]\label{iq:s1:short-term-reversal:2}
\omterm{def:s1:short-term-reversal:residual}{Retournement résiduel} face au retournement brut~: quelle est la différence, et pourquoi est-elle importante~?
\end{interviewq}

\begin{interviewq}[$\star\star$ \iqroles{trader}]\label{iq:s1:short-term-reversal:3}
Une stratégie de retournement affiche un ratio de Sharpe brut de 7 et un coût de point mort de 15 points de base. La mettriez-vous en œuvre~? Que demanderiez-vous~?
\end{interviewq}

\begin{interviewq}[$\star\star$ \iqroles{développeur, chercheur}]\label{iq:s1:short-term-reversal:4}
Comment construiriez-vous un optimiseur quotidien tenant compte des coûts pour un millier de titres, qui tourne en quelques secondes~?
\end{interviewq}

\begin{interviewq}[$\star\star$ \iqroles{risque}]\label{iq:s1:short-term-reversal:5}
Quand un portefeuille de retournement perd-il rapidement de l'argent, et que surveilleriez-vous~?
\end{interviewq}

\begin{interviewq}[$\star\star\star$ \iqroles{chercheur}]\label{iq:s1:short-term-reversal:6}
Sous un impact en racine carrée, dérivez la taille de l'ordre à partir de laquelle une prévision $\alpha$ à un jour couvre tout juste son coût, pour un titre de volatilité $\sigma$, de volume quotidien
$V$ et de demi-fourchette $h$.
\end{interviewq}
